\chapter*{Introduction} % chapter* je necislovana kapitola
\addcontentsline{toc}{chapter}{Introduction} % rucne pridanie do obsahu
\markboth{INTRODUCTION}{INTRODUCTION} % vyriesenie hlaviciek

A \emph{flow} on a graph is an assignment of flow values to the darts of the graph. For convenience, each edge is replaced by two darts, and when we use a flow value on some edge, we can choose the dart in the desired direction. The flow must satisfy two conditions -- flow values on opposite darts are additive inverses and the sum of flow values surrounding a vertex is always zero. When the flow value of some edge is zero, then the given flow is de facto a flow on the graph after deleting that edge. Hence, we usually require the flow to be \emph{nowhere-zero}. Moreover, we can restrict the range of flow values permitted. A \emph{nowhere-zero $k$-flow} can use only integral values between $1$ and $k-1$ and their additive inverses. The smallest number $k$ such that a given graph admits a $k$-flow is called a flow number.

Research on integral nowhere-zero flows started in the early $1950$s by Tutte \cite{tutte_beginning}. This introductory paper was followed by another paper of Tutte, in which he described the relationship between the nowhere-zero flows and vertex colourings of graphs and conjectured that each bridgeless graph admits a nowhere-zero $5$-flow \cite{tutte}.

The main problem in this area is finding the best upper bound on the flow number. Since there are graphs that do not admit a $4$-flow, the bound cannot be better than $5$. One of the first proved bounds was the one of Jaeger \cite{jaeger1979flows}, when he proved that every bridgeless graph has a nowhere-zero $8$-flow. Seymour \cite{seymour} improved this result, showing that each bridgeless graph allows a nowhere-zero $6$-flow. Note that the gap between the known and conjectured results was narrowed to only two possibilities -- either there exists a nowhere-zero $5$-flow on any bridgeless graph, or there is a counterexample that necessarily requires a flow value $5$ or $-5$.

There are several properties proved about the minimal counterexample to the $5$-flow conjecture (if there exists any) -- it must be cubic, not $3$-edge-colourable, cyclically $6$-edge-connected \cite{min_ctrex_cyc_edge_con}, of girth at least $12$ \cite{min_ctrex_girth} and oddness at least $6$ \cite{oddness_4}. However, it is not even known whether a graph with such properties exists.

The possible counterexample to the flow problem needs to be very large. The smallest cubic graph of girth at least $12$, the Benson graph, has $126$ vertices \cite{benson}. Moreover, this graph is colourable. Hence, it is not expected that we would be able to find the counterexample by an exhaustive generation of all snarks, since the currently fastest generators \texttt{tetration} and \texttt{minisnark} \cite{tetration_minisnark} are able to generate snarks up to $40$ vertices.

However, also the local structure of the graph may be interesting in the investigation of the problem, which motivates the generation of cubic multipoles. Although there exist some generators (mainly the \texttt{pregraphs} \cite{pregraph} and the unpublished \texttt{multigraph}), they are not fast enough and the \texttt{pregraphs} has no option to prescribe minimal girth (which is obviously the crucial part). Therefore, we develop a generator \texttt{polehunter}, for now able only to generate mutlipoles of minimal girth up to $5$, which is faster than the currently known generators.

In the meantime, people tried to relax the constraints of the flow problem. One of the first ideas was to use fractional flow values \cite{goddyn_fractional_nzf}. There were also attempts to use real flow values, but Steffen showed that this relaxation does not help \cite{steffen_circular}.

Another relaxation assigns multidimensional flow values. This idea was brought to light by Jain \cite{3d_unit_flows}, who conjectured that in three dimensions, Euclidean unit vectors are sufficient to construct a flow on any bridgeless graph. Thomassen \cite{7d_unit_flows} proved this in seven dimensions. However, these papers focused only on the unit vectors and did not generalise the concept to the graphs, which do not admit a flow using only unit vectors.

Since in three dimensions, it is already conjectured that all graphs admit the unit vector flow, the first paper concerning multidimensional flows not using only unit vectors focused on the two-dimensional case \cite{euclidean_flows}. Considering a $2$-dimensional $r$-flow, the bounding criterion on the permitted flow values is their Euclidean norm, which must be at least $1$ and at most $r-1$. In this setting, the flow number is often irrational. That indicates that the choice of the norm is not ideal.

Therefore, in recent preprints, the idea of Manhattan and Chebyshev flows appeared \cite{preprint, p-norm_flows}. With respect to these norms, the flow number is always rational. The proved and the conjectured upper bounds are $3$ and $5/2$, respectively. The unit vectors are proved to be sufficient also in seven dimensions w.r.t. the Manhattan norm and in three dimensions w.r.t. the Chebyshev norm. Again, the $2$-dimensional problem is interesting for snarks, since for and only for colourable cubic graphs, a flow composed of unit vectors always exists. The research also led to a pair of flows, where the first one uses values from $\{0, \pm 1\}$, the other values $\{0, \pm 1, \pm 2, \pm 3\}$, but moreover, $(0,0), (0,1)$ and $(0,-1)$ are forbidden. This so-called \emph{flow-pair} implies the existence of $5$-flow and two-dimensional Chebyshev $5/2$-flow.

In this work, we show other interesting relations of this flow-pair. We prove that the flow-pair exists on snarks of oddness $2$. For graphs with a flow-pair, we are also able to construct a nowhere-zero $7$-flow not using the same element around a vertex twice (called a \emph{rich} flow). Thus, in total, graphs of oddness $2$ admit a rich nowhere-zero $7$-flow.

As mentioned above, snarks are exactly those cubic graphs, for which flow number is greater than two. One interesting question is, how small it could be. We managed to prove a lower bound dependent on the order of a graph, already submitted to a journal \cite{preprint}, and similar to a one-dimensional lower bound on flow number \cite{cycle_rank}. We also use a generalised version of the flow-pair to construct an infinite class of snarks with flow numbers, that can be arbitrarily small.